\documentclass{article}
\usepackage[utf8]{inputenc}
\usepackage[T1]{fontenc}
\usepackage[margin=2cm]{geometry}
\usepackage{hyperref}
\sloppy
\begin{document}
\section*{Historical source}
Encyclopaedia Britannica scan, source row 1939. Parse 3.5 extracted the text. Super 3.5 VL corrected it from an enlarged crop of each block. It can still contain recognition errors.
Page headers, footers, and page numbers are removed.

\subsection*{PHARMACY (E0169)}
PHARMACY.

Part II.

Preparations.

flower. When in perfection, it is very light and thin; pale, and almost colourless; of great fragrance, tho' not quite so agreeable as the remedy itself. It is recommended, in the dose of a few drops, in nervous and hysteric complaints. Boerhaave holds it in great esteem against epilepsies, and suppressions of the uterine purgations, occasioned by weakness and inactivity.

535. Essential oil of rue-leaves. L.

The oil of rue has a very acrid taste, and a penetrating smell, resembling that of the herb, but rather more unpleasant. It is sometimes made use of in hysteric disorders, and as an anthelmintic; as also in epilepsies proceeding from a relaxed state of the nerves.

Rue yields its oil very sparingly. The largest quantity is obtained from it when the flowers are ready to fall off, and the seeds begin to flow themselves: suitable maceration, previous to the distillation, is here extremely necessary.

536. Essential oil of savin leaves. L. E.

This oil is a celebrated uterine and emmenagogue ; in cold phlegmatic habits, it is undoubtedly a medicine of good service, though not capable of performing what it has been usually represented to do. The dose is, two or three drops or more.

537. Essential oil of saffron. L. E.

This is the most ponderous of all the known essential oils, but rises in distillation with sufficient ease: it appears limpid as water, has a moderately pungent taste, a very fragrant smell, exactly resembling that of the laffafras. It finds greatly commended as a sudorific, and for purifying the blood and juices: it is likewise supposed to be of service in humourous asthmas and coughs. The dose is from one drop to eight or ten; though Geoffroy goes as far as 20.

The decoction remaining after the distillation of the oil, affords, by inspissation, an useful extract, of a mild, bitterish, subastringent taste. Hoffman says, he has given it with great benefit, in doses of a scruple, as a corroborant in cachectic cases, in the decline of intermitting fevers, and for abating hypochondriacal spasms.

538. Oil of turpentine. L. E.

This is distilled in the same manner as the foregoing oils, and is strictly an essential one, though not usually ranked in this class: it is commonly, but improperly, called spirit of turpentine. It is employed in large quantities for some mechanic purposes, and hence the distillation of it is become a particular business.

This oil is a very hot, stimulating medicine. It is sometimes given as a sudorific and diuretic, in the dose of two or three drops: in larger doses, it is apt to greatly heat the body, occasion pain of the head, an effusion of the semen and liquor of the prostate glands. It has nevertheless been of late taken in considerable doses (along with honey or other convenient vehicles) against the sciatica; and, as is said, with good success. Some have recommended it against venereal runnings: but here it has produced mischievous consequences, inflaming the parts and aggravating the disorder. Externally it is not unfre-

uently employed against rheumatic pains, aches, sprains, for difcuffing cold tumours, and restraining hemor- rhages.

539. After the distillation of the turpentine, there remains in the still a brittle resinous substance, of a yellow colour, called yellow resin [L.]

The only use of this is in external applications, for giving consistence to platters, and the like purpofes.

540. Most of the foregoing oils are drawn by our chemists, and easily procurable in a tolerable degree of perfection; those of cinnamon, cloves, nutmegs, and mace, excepted. These are usually imported from abroad; and are for the most part so much adulterated, that it is difficult to meet with such as are at all fit for use.

Nor are the adulterations of these kinds of prepara-tions easily discoverable. The grosser abufes indeed may be readily detected: thus if the oil is mixed with spirit of wine, it will turn milky on the addition of water; if with expressed oils, rectified spirit will dissolve the efential, and leave the other behind; if with oil of turpentine, on dipping a piece of paper in the mixture, and drying it with a gentle heat, the turpen-tine will be betrayed by its fmel. But the more subtile artifis have contrived other methods of fophifcating, which elude all trials of this kind.

Some have looked upon the specific gravity of oils, as a certain criterion of their genuineness. This, however, is not to be absolutely depended on; for the genuine oils, obtained from the same subject, often-times differ in gravity as much as those drawn from different ones. Cinnamon and cloves, whose oils usually sink in water, yield, if slowly and variably distilled, an oil of great fragrancy, which is nevertheless specifically lighter than the aqueous fluid employed in the distillation of it; whilst, on the other hand, the last runnings of some of the lighter oils prove sometimes so ponderous as to sink in water.

As all essential oils agree in the general properties of solubility in spirit of wine, indissolubility in water, miscibility with water by the intervention of certain intermedia, volatility in the heat of boiling water, \&c. it is plain that they may be variously mixed with one another, or the dearer sophisticated with the cheaper, without any possibility of discovering the abuse by any trials of this kind. And indeed, it would not be of much advantage to the purchaser if he had infallible criteria of the genuineness of every individual oil. It is of as much importance, that they may be good, as that they may be genuine; for we have often seen genuine oils, from incautious distillation, and long and careful keeping, weaker both in smell and taste than the common sophisticated ones.

The smell and taste seem to be the only certain tests that the nature of the thing will admit of. If a bark should have in every respect the appearance of good cinnamon, and should be proved indubitably to be the genuine bark of the cinnamon tree; yet, if it wants the cinnamon flavour, or has it but in a low degree, we reject it: and the case is the same with the oil. It is only from use and habit, or comparison with specimens of known quality, that we can judge of the goodness, either of the drugs themselves, or of their oils.


\section*{Source record}
\begin{description}
\item[Scan ID] \texttt{britannica:309f765b1ecc83e0a7f7762e44b4a927a7555443:train:1939}
\item[Image file] \texttt{page-1939-0f50bb0b2f40.jpg}
\item[Image SHA-256] \texttt{0f50bb0b2f406428516bb469f72abc9fc9fa54cdb87fbdff0adae232ef9c2384}
\item[Dataset] NationalLibraryOfScotland/encyclopaedia\_britannica\_illustrated at revision \texttt{309f765b1ecc83e0a7f7762e44b4a927a7555443} (CC0-1.0)
\item[Extraction] model \texttt{nvidia/nemotron-3.5-parse}, raw reply \texttt{extract/raw/page-1939-0f50bb0b2f40.json}
\item[Raw reply SHA-256] \texttt{a055b80947edd215e5844732cb75fbf7dad181cfef63a25e4992d434abb28af3}
\item[Correction] model \texttt{nvidia/Nemotron-3.5-Super-VL}, block crops, raw replies \texttt{correct/raw/crop\_nothink/page-1939-0f50bb0b2f40-b*.json}; super\_vl: 23
\item[Running header] none
\end{description}
\end{document}
